% GENERATED FILE. Edit canonical lessons, then run npm run generate:books.
% canonical-source-hash: fnv1a64:97a8f148297ebc55
% canonical-lessons: ZH-C300-dizhen, ZH-C300-si, ZH-W300-si-die, ZH-C300-shengming, ZH-C300-siyin, ZH-C300-baozha

\chapter{Life and Death}
\label{ch:zh-life-and-death}
\hlchaptermodality{300}

\begin{chapteropening}
\textbf{By the end of this chapter:} \emph{I can talk about an earthquake and an explosion, who died, and what they died of.}

Complete the last lesson of chapter 300: I can talk about an earthquake and an explosion, who died, and what they died of.
\end{chapteropening}

\section[dìzhèn]{dìzhèn — earthquake}

\label{lesson:ZH-C300-dizhen}



\begin{quote}
Before the new one: say the Mandarin for to pass out, then the Mandarin for to fight to save a patient.
\end{quote}

\subsection*{You'll want to know: dìzhèn}
\emph{dìzhèn} — \textquotedblleft{}earthquake\textquotedblright{}.

\begin{quote}
\emph{zhèr yǐqián yǒu guo dìzhèn} — There was once an earthquake here.
\end{quote}

Its characters come later; for now it is a word for the ear.

\begin{cousinweb}
\emph{dìzhèn} is built from \zh{地} 'the ground' and \emph{zhèn} 'to shake', both inherited from Middle Chinese, \emph{dijH} and \emph{tsyinH} (Guangyun, in Baxter's transcription), which glosses the second 'thunder' and 'to move': the ground shaking.
\end{cousinweb}

\begin{grammarlens}[title={What you've built}]
One more word for telling what happened.
\end{grammarlens}

\subsection*{Your turn}
\begin{itemize}
\raggedright
  \item \emph{Say it:} \emph{dìzhèn}
  \item \emph{Say it:} \emph{dìzhèn}, once more
  \item \emph{Say it:} \emph{zhèr yǐqián yǒu guo dìzhèn} — There was once an earthquake here.
  \item \emph{Recall:} say the Mandarin for to pass out, then the Mandarin for to fight to save a patient, then say \emph{dìzhèn} again
\end{itemize}

\begin{tcolorbox}[breakable,colback=teal!4,colframe=teal!35!black,title={Before you move on}]
What does dìzhèn mean? (\textquotedblleft{}Earthquake\textquotedblright{}.) Say it once more.
\end{tcolorbox}
\section[\texorpdfstring{\zh{死} (sǐ)}{死 (sǐ)}]{\zh{死} (sǐ) — to die}

\label{lesson:ZH-C300-si}



\begin{quote}
Before the new one: say the Mandarin for to fight to save a patient, then the Mandarin for an earthquake.
\end{quote}

\subsection*{You'll want to know: \zh{死}}
\textbf{\zh{死}} — \emph{sǐ} — \textquotedblleft{}to die\textquotedblright{}.

\begin{quote}
\textbf{\zh{我的花都死了}} \emph{wǒ de huā dōu sǐ le} — All my flowers have died.
\end{quote}

The new character, in pieces: \zh{死} — six strokes: a long bar on top; under it on the left the evening piece of \zh{外}, three strokes; on the right a short sweep and a vertical that bends right along the foot and hooks up, as on the right of \zh{比}. The next lesson puts it in your hand.

\begin{cousinweb}
\textbf{\zh{死}} (\emph{sǐ}) is inherited from Middle Chinese \emph{sijX} (Guangyun, in Baxter's transcription), which quotes an older dictionary for it: 'to be used up', what a person parts from.
\end{cousinweb}

\begin{grammarlens}[title={What you've built}]
One more word for telling what happened.
\end{grammarlens}

\subsection*{Your turn}
\begin{itemize}
\raggedright
  \item \emph{Say it:} \emph{sǐ}
  \item \emph{Say it:} \emph{sǐ}, once more
  \item \emph{Say it:} \emph{wǒ de huā dōu sǐ le} — All my flowers have died.
  \item \emph{Recall:} say the Mandarin for to fight to save a patient, then the Mandarin for an earthquake, then say \emph{sǐ} again
\end{itemize}

\begin{tcolorbox}[breakable,colback=teal!4,colframe=teal!35!black,title={Before you move on}]
What does \zh{死} mean? (\textquotedblleft{}To die\textquotedblright{}.) Say it once more.
\end{tcolorbox}
\section[\texorpdfstring{\zh{死} (sǐ)}{死 (sǐ)}]{\zh{死} — a long bar over two small pieces}

\label{lesson:ZH-W300-si-die}



\begin{quote}
The word first, then the character inside it.

\begin{itemize}
\raggedright
  \item \emph{Recall:} say \emph{sǐ} — \textbf{R1}, one lesson back, before it decays
\end{itemize}

\begin{quote}
\zh{外} \zh{比} \zh{多}
\end{quote}

\emph{Read it:} each one aloud, with nothing said first — \emph{wài}, \emph{bǐ}, \emph{duō}
\end{quote}

\subsection*{Script — a long bar over two small pieces}
\hlblockfigure{\detokenize{figures/ZH-W300-si-die-filmstrip.pdf}}{How \zh{死} is written, stroke by stroke}

\begin{quote}
\zh{死}
\end{quote}

\textbf{Three pieces: a bar on top, then two side by side under it.}

\textbf{The bar, one stroke}, long, across the top.

\textbf{Under it on the left, three strokes.} The evening piece on the left of \zh{外}: a short sweep, a bar that turns into a long sweep to the lower left, and a dot.

\textbf{On the right, two strokes.} The right half of \zh{比}: a short sweep down to the left, then a vertical that bends right along the foot and hooks up.

1 + 3 + 2 = six.

\subsection*{Your turn}
\begin{itemize}
\raggedright
  \item \emph{Write it:} \zh{死}, all six, with the model in front of you
  \item \emph{Cover:} write it again from memory, and count the pieces, not the strokes
\end{itemize}

\begin{tcolorbox}[breakable,colback=teal!4,colframe=teal!35!black,title={Before you move on}]
How many strokes in \zh{死}? (\textbf{Six}.) What comes first? (\textbf{The long bar on top}.) Where have you made the right piece? (\textbf{On the right of \zh{比}}.)

Source: \href{https://raw.githubusercontent.com/chanind/hanzi-writer-data/68d10a4b21150cae5e1ebbd223eed289cf32d90c/data/\%E6\%AD\%BB.json}{Hanzi Writer Data \zh{死}, pinned stroke paths and medians}.
\end{tcolorbox}
\section[shēngmìng]{shēngmìng — the life of a living thing}

\label{lesson:ZH-C300-shengming}



\begin{quote}
Before the new one: say the Mandarin for an earthquake, then the Mandarin for to die.
\end{quote}

\subsection*{You'll want to know: shēngmìng}
\emph{shēngmìng} — \textquotedblleft{}the life of a living thing\textquotedblright{}.

\begin{quote}
\emph{měi ge rén de shēngmìng dōu hěn zhòngyào} — Every person's life matters.
\end{quote}

Its characters come later; for now it is a word for the ear.

\begin{cousinweb}
\emph{shēngmìng} is built from \zh{生} 'to live' and \emph{mìng}, both inherited from Middle Chinese, \emph{srjaeng} and \emph{mjaengH} (Guangyun, in Baxter's transcription), which glosses the first 'to grow' and the second 'to order'; its sense 'a life' is not given there.
\end{cousinweb}

\begin{grammarlens}[title={What you've built}]
One more word for telling what happened.
\end{grammarlens}

\subsection*{Your turn}
\begin{itemize}
\raggedright
  \item \emph{Say it:} \emph{shēngmìng}
  \item \emph{Say it:} \emph{shēngmìng}, once more
  \item \emph{Say it:} \emph{měi ge rén de shēngmìng dōu hěn zhòngyào} — Every person's life matters.
  \item \emph{Recall:} say the Mandarin for an earthquake, then the Mandarin for to die, then say \emph{shēngmìng} again
\end{itemize}

\begin{tcolorbox}[breakable,colback=teal!4,colframe=teal!35!black,title={Before you move on}]
What does shēngmìng mean? (\textquotedblleft{}The life of a living thing\textquotedblright{}.) Say it once more.
\end{tcolorbox}
\section[\texorpdfstring{\zh{死因} (sǐyīn)}{死因 (sǐyīn)}]{\zh{死因} (sǐyīn) — cause of death}

\label{lesson:ZH-C300-siyin}



\begin{quote}
Before the new one: say the Mandarin for to die, then the Mandarin for the life of a living thing.
\end{quote}

\subsection*{You'll want to know: \zh{死因}}
\textbf{\zh{死因}} — \emph{sǐyīn} — \textquotedblleft{}cause of death\textquotedblright{}.

\begin{quote}
\textbf{\zh{他的死因是什么}} \emph{tā de sǐyīn shì shénme} — What did he die of?
\end{quote}

\begin{cousinweb}
\textbf{\zh{死因}} (\emph{sǐyīn}) is built from \zh{死} 'to die' and \zh{因} 'a cause', both inherited from Middle Chinese, \emph{sijX} and \emph{'jin} (Guangyun, in Baxter's transcription), which glosses the second 'to rely on' and 'a reason'.
\end{cousinweb}

\begin{grammarlens}[title={What you've built}]
One more word for telling what happened.
\end{grammarlens}

\subsection*{Your turn}
\begin{itemize}
\raggedright
  \item \emph{Say it:} \emph{sǐyīn}
  \item \emph{Say it:} \emph{sǐyīn}, once more
  \item \emph{Say it:} \emph{tā de sǐyīn shì shénme} — What did he die of?
  \item \emph{Recall:} say the Mandarin for to die, then the Mandarin for the life of a living thing, then say \emph{sǐyīn} again
\end{itemize}

\begin{tcolorbox}[breakable,colback=teal!4,colframe=teal!35!black,title={Before you move on}]
What does \zh{死因} mean? (\textquotedblleft{}Cause of death\textquotedblright{}.) Say it once more.
\end{tcolorbox}
\section[bàozhà]{bàozhà — to explode; to blow up}

\label{lesson:ZH-C300-baozha}



\begin{quote}
Before the new one: say the Mandarin for the life of a living thing, then the Mandarin for the cause of death.
\end{quote}

\subsection*{You'll want to know: bàozhà}
\emph{bàozhà} — \textquotedblleft{}to explode; to blow up\textquotedblright{}.

\begin{quote}
\emph{zuótiān wǎnshang fāshēng le bàozhà} — There was an explosion last night.
\end{quote}

Its characters come later; for now it is a word for the ear.

\begin{cousinweb}
\emph{bàozhà} is built from \emph{bào} 'to burst' and \emph{zhà} 'to blow apart'. The first is inherited from Middle Chinese \emph{paewH} (Guangyun, in Baxter's transcription), which glosses it 'fire cracking'; the Guangyun has no entry for the second, so its history is not settled.
\end{cousinweb}

\begin{grammarlens}[title={What you've built}]
One more word for telling what happened.
\end{grammarlens}

\subsection*{Your turn}
\begin{itemize}
\raggedright
  \item \emph{Say it:} \emph{bàozhà}
  \item \emph{Say it:} \emph{bàozhà}, once more
  \item \emph{Say it:} \emph{zuótiān wǎnshang fāshēng le bàozhà} — There was an explosion last night.
  \item \emph{Recall:} say the Mandarin for the life of a living thing, then the Mandarin for the cause of death, then say \emph{bàozhà} again
\end{itemize}

\begin{tcolorbox}[breakable,colback=teal!4,colframe=teal!35!black,title={Before you move on}]
What does bàozhà mean? (\textquotedblleft{}To explode; to blow up\textquotedblright{}.) Say it once more.
\end{tcolorbox}
